\documentclass[11pt,a4paper]{article}
% Shared preamble of the FEBAUTO Racing manuals.
\usepackage[margin=2.2cm,headheight=16pt]{geometry}
\usepackage[T1]{fontenc}
\usepackage{lmodern}
\usepackage{xcolor}
\usepackage{graphicx}
\usepackage{booktabs}
\usepackage{tabularx}
\usepackage{listings}
\usepackage{fancyhdr}
\usepackage{enumitem}
\usepackage{tikz}
\usetikzlibrary{arrows.meta}
\usepackage[hidelinks,colorlinks=true,linkcolor=febblue,urlcolor=febblue]{hyperref}

\definecolor{febblue}{HTML}{1C4687}
\definecolor{febgold}{HTML}{FDB515}
\definecolor{codebg}{HTML}{F3F4F6}

\lstset{
  basicstyle=\ttfamily\small,
  backgroundcolor=\color{codebg},
  frame=single, framerule=0pt, framesep=4pt,
  xleftmargin=4pt, xrightmargin=4pt,
  breaklines=true, columns=fullflexible, keepspaces=true,
  showstringspaces=false,
}
\newcommand{\file}[1]{\texttt{#1}}
\newcommand{\cmd}[1]{\texttt{#1}}
\newcommand{\rulebox}[1]{\par\noindent\colorbox{codebg}{\parbox{\dimexpr\linewidth-2\fboxsep}{#1}}\par\medskip}

\pagestyle{fancy}
\fancyhf{}
\fancyhead[L]{\color{febblue}\textbf{FEBAUTO Racing}}
\fancyhead[R]{\color{febblue}\nouppercase{\leftmark}}
\fancyfoot[C]{\thepage}
\renewcommand{\headrulewidth}{0.8pt}
\renewcommand{\headrule}{\color{febgold}\hrule width\headwidth height\headrulewidth}
\setcounter{tocdepth}{1}

% \febtitle{subtitle}{part lines}
\newcommand{\febtitle}[2]{%
\begin{titlepage}
  \centering
  \vspace*{2cm}
  \includegraphics[width=5cm]{../../assets/feb_logo_512.png}\\[1.5cm]
  {\color{febblue}\Huge\bfseries FEBAUTO Racing}\\[6pt]
  {\Large #1}\\[2cm]
  \begin{tabular}{rl}#2\end{tabular}\\[3cm]
  Formula Electric at Berkeley\\
  September 2026, platform release v0.1.0\\[6pt]
  \url{https://github.com/FEBAutonomous/feb-racing}\\
  \url{https://febautonomous.github.io/feb-racing/}
\end{titlepage}
\tableofcontents
\newpage}

\setlist{nosep,leftmargin=*}

\begin{document}
\febtitle{Release 2: keeping your driver}{%
  \textbf{For} & anyone who wrote code in the old devkit's \file{stack/} folder\\[3pt]
  \textbf{Time} & about 30 minutes, 20 of them a download\\[3pt]
  \textbf{Result} & your driver running in the new box, ready to score and submit\\}

\section{Why you cannot just pull}

Release 2 rewrote the repository's history. In an old clone, \cmd{git pull} stops with
\emph{refusing to merge unrelated histories}, and the old launcher \cmd{feb-sim} no longer exists.
You clone once more, and carry your \file{stack/} folder across by hand. Nothing else in the old
clone is yours. After this, \cmd{./feb update} covers every future release; you never do this again.

\section{Save your work}

Your code is in the old clone's \file{stack/} folder, and it is not on GitHub. Copy it out
\emph{before} anything else:

\begin{lstlisting}
cp -r ~/feb-racing/stack ~/feb-stack-backup
ls ~/feb-stack-backup
\end{lstlisting}

You should see your package folders. If you changed anything outside \file{stack/}, copy that
too. Then delete the old clone:

\begin{lstlisting}
rm -rf ~/feb-racing
\end{lstlisting}

Windows: do all of this in your WSL 2 Ubuntu terminal.

\section{Set up release 2}

If Docker, VS Code and Foxglove are already installed, this is the whole setup:

\begin{lstlisting}
mkdir -p ~/FEB && cd ~/FEB
git clone git@github.com:FEBAutonomous/feb-racing.git
cd feb-racing
./feb setup --name "Your Name"
\end{lstlisting}

It pulls the new image, downloads the simulator and creates your box, a persistent container with
your whole \file{\textasciitilde/FEB} folder inside. Then \cmd{./feb up}: the simulator opens, arrow keys
drive. Close the window to stop. That proves the new devkit works before your code goes in.

\section{Put your driver back}

\begin{lstlisting}
cp -r ~/feb-stack-backup/* ~/FEB/feb-racing/stack/
ls ~/FEB/feb-racing/stack
\end{lstlisting}

Your packages are now in the new \file{stack/}, which is the workspace the box builds.

\section{What changed for your code}

Three things differ from release 1, and the first two matter for every driver.

\subsection*{1. The topics: the car's, not the simulator's}

Release 1 drivers subscribed to the simulator directly, \cmd{/autodrive/roboracer\_1/lidar} and
published \cmd{/autodrive/roboracer\_1/throttle\_command}. Release 2 drivers talk the car's topics,
and a house node (the \emph{sim adapter}) translates them for the simulator. The same code then runs
on the car unchanged.

\begin{tabularx}{\linewidth}{lll}
\toprule
Was & Now & Type \\
\midrule
\cmd{/autodrive/roboracer\_1/lidar} & \cmd{/scan} & sensor\_msgs/LaserScan \\
\cmd{/autodrive/roboracer\_1/imu} & \cmd{/imu} & sensor\_msgs/Imu \\
encoders, ips & \cmd{/odom} & nav\_msgs/Odometry (dead-reckoned) \\
\cmd{/autodrive/roboracer\_1/front\_camera} & \cmd{/camera/color/image\_raw} & sensor\_msgs/Image \\
\cmd{throttle\_command} + \cmd{steering\_command} & \cmd{/drive} & ackermann\_msgs/AckermannDriveStamped \\
\bottomrule
\end{tabularx}

\medskip
On \cmd{/drive} you publish a \textbf{speed in m/s} and a \textbf{steering angle in radians}, left
positive. No more throttle fractions. In the gap follower this is about thirty lines; the starter
kit's \file{starter\_kit/feb\_driver/feb\_driver/reactive\_driver.py} shows the new shape.

\subsection*{2. The driver's card}

Each driver package now carries a \file{driver.yaml} at its top level, so \cmd{feb up}, \cmd{feb
practice} and the car know how to start it:

\begin{lstlisting}
name: my_driver
kind: walls                 # walls | cones | any
maintainer: your name
packages: [my_driver]       # what feb submit copies: your own packages
modes:
  trackdrive: {launch: drive.launch.py, max_speed: 3.0}
  inspection: {launch: drive.launch.py, max_speed: 0.8}
default_mode: trackdrive
\end{lstlisting}

\cmd{launch} is a file in your package's \file{launch/} folder. If you had no launch file, copy
\file{starter\_kit/feb\_driver/launch/drive.launch.py} and change the package and executable names.

\subsection*{3. Config files (optional)}

Tuning numbers go in \file{config/driver.yaml}, loaded by the launch file; values that differ on
the real car go in \file{config/driver.car.yaml}, loaded on top only on the car. Not required to run.

\section{Run it}

\begin{lstlisting}
./feb up --driver my_driver
\end{lstlisting}

The box builds your package with \cmd{colcon} on the first run, then every run after that is
incremental. Build errors print in this terminal. Until your driver has its card, you can still run
it the raw way:

\begin{lstlisting}
./feb up --stack "ros2 launch my_driver drive.launch.py"
\end{lstlisting}

\section{Look inside}

\begin{lstlisting}
./feb shell              # a terminal in the box: ros2 topic list, ros2 topic echo /drive
./feb shell --gui        # the same, plus the box's desktop in a browser tab (localhost:6080)
\end{lstlisting}

\textbf{rviz in the browser:} with \cmd{./feb up} running, in a second terminal \cmd{./feb shell -{}-gui}
then \cmd{rviz2 \&}. On the browser desktop: Fixed Frame \cmd{odom}, Add by topic \cmd{/scan} and
\cmd{/odom}, LaserScan style Points, size 4; Reliability Best Effort if a topic stays empty.
\cmd{./feb code} opens VS Code inside the box.

\section{Score and submit}

\begin{lstlisting}
./feb practice --driver my_driver      # ten laps on each qualifying track
./feb submit                           # results to the scoreboard, code to feb-stack
\end{lstlisting}

A run counts only with zero collisions. Everything else, the dead-man, the car, is in
\file{member\_guide.pdf}.

\section{When something is wrong}

\begin{tabularx}{\linewidth}{lX}
\toprule
It says & Do \\
\midrule
\texttt{no driver.yaml with name 'my\_driver'} & no card with that \cmd{name:} under \file{stack/}; add it (section 5.2) \\
\texttt{refusing to merge unrelated histories} & you are in the old clone; follow this document from the top \\
nothing drives, \cmd{/drive} is silent & your node still publishes the old throttle topics; see ``The topics'' \\
build error mentioning \texttt{autodrive\_msgs} & that package is gone; the adapter owns the simulator topics now \\
\texttt{this clone sits directly in /home/...} & \cmd{mkdir \textasciitilde/FEB \&\& mv \textasciitilde/feb-racing \textasciitilde/FEB/} \\
\texttt{can't copy '...resource/my\_driver'} & you renamed the package but not the empty marker file in \file{resource/}; rename it to the package name \\
\bottomrule
\end{tabularx}

\end{document}
